\documentclass[11pt]{letter}
\usepackage[T1]{fontenc}
\usepackage[utf8]{inputenc}
\usepackage{lmodern}
\usepackage[margin=1in]{geometry}
\usepackage{microtype}

\signature{Darcio Genicolo-Martins\\INSPER, S\~ao Paulo\\\texttt{darciogm1@insper.edu.br}}
\address{Darcio Genicolo-Martins\\INSPER\\R.\ Quat\'a 300, S\~ao Paulo, SP, 04546-042, Brazil}

\begin{document}

\begin{letter}{The Editors\\\emph{Journal of Public Economics}}

\opening{Dear Editors,}

Please find attached my manuscript, \emph{Sheltered Bidding: The Within-Auction Cost of SME Set-Asides}, for consideration at the \emph{Journal of Public Economics}.

The paper studies the welfare cost of SME set-asides in public procurement using a March~2018 legal reversal on São Paulo's centralized procurement platform. A difference-in-differences design shows that extending SME-only procurement to medical supplies raised winning prices by roughly 10~percent and doubled SME participation. I then use the platform's Pregão drop-out format to estimate an asymmetric IPV auction model and decompose the price effect into a within-auction order-statistic component and an endogenous-entry component.

The main finding is that two-thirds to three-quarters of the simulated price effect comes from the within-auction component: bidder exclusion changes which order statistic clears the restricted auction. Endogenous SME entry partially offsets, but does not eliminate, this cost. The implied welfare loss reaches R\$\,55~million per year on a single product group, and a 10~percent price preference welfare-dominates set-asides in thick standardized markets, while the ranking becomes conditional in thinner pharmaceutical markets.

The paper contributes to the public economics of procurement preferences by connecting reduced-form evidence, auction-based structural decomposition, and welfare comparison across alternative policy instruments.

Thank you for your consideration.

\closing{Sincerely,}

\end{letter}

\end{document}
